\section{Strings}

\subsection{KMP}

Encontra a primeira ocorrência de um padrão em uma string
em $O(N+M)$, onde $N$ é o tamanho do texto e $M$ é o tamanho
do padrão.

\textbf{Prefix Function:} $O(M)$.

\textbf{Busca:} $O(N)$.

A prefix function armazena, para cada posição, o tamanho do maior
prefixo próprio que tambem é um sufixo.

\begin{lstlisting}
vector<int> prefixFunction(const string& pattern) {
    int patternLength = pattern.size();
    vector<int> prefix(patternLength);

    for (int i = 1; i < patternLength; i++) {
        int matchedLength = prefix[i - 1];

        // Recuar para o maior prefixo que ainda pode ser aproveitado.
        while (matchedLength > 0 && pattern[i] != pattern[matchedLength]) 
            matchedLength = prefix[matchedLength - 1];

        if (pattern[i] == pattern[matchedLength])
            matchedLength++;

        prefix[i] = matchedLength;
    }

    return prefix;
}

int kmp(const string& text, const string& pattern) {
    vector<int> prefix = prefixFunction(pattern);

    int patternLength = pattern.size();

    for (int textIndex = 0, matchedLength = 0; textIndex < text.size(); textIndex++) {
        // Usa a prefix function para evitar voltar no texto.
        while (matchedLength > 0 && text[textIndex] != pattern[matchedLength])
            matchedLength = prefix[matchedLength - 1];

        if (text[textIndex] == pattern[matchedLength])
            matchedLength++;

        // O padrao terminou neste indice do texto.
        if (matchedLength == patternLength)
            return textIndex - patternLength + 1;
    }

    return -1;
}
\end{lstlisting}

Uso:
\begin{lstlisting}
string text, pattern;
cin >> text >> pattern;

int position = kmp(text, pattern);

cout << position << '\n';
\end{lstlisting}

Para índice baseado em $1$:

\begin{lstlisting}
cout << kmp(text, pattern) + 1 << '\n';
\end{lstlisting}


\subsection{Z-Function}

Calcula, para cada posição, o tamanho do maior prefixo da string
que coincide com a substring iniciada naquela posição.

\textbf{Complexidade:} $O(N)$.

\begin{lstlisting}
vector<int> zFunction(const string& text) {
    int n = text.size();
    vector<int> z(n);

    int left = 0;
    int right = 0;

    for (int i = 1; i < n; i++) {
        // Aproveita uma correspondencia ja calculada.
        if (i <= right)
            z[i] = min(right - i + 1, z[i - left]);

        // Expande enquanto os caracteres continuarem iguais.
        while (i + z[i] < n && text[z[i]] == text[i + z[i]]) {
            z[i]++;
        }

        // Atualiza o intervalo com a maior correspondencia atual.
        if (i + z[i] - 1 > right) {
            left = i;
            right = i + z[i] - 1;
        }
    }

    return z;
}
\end{lstlisting}

Uso:
\begin{lstlisting}
string text;
cin >> text;

vector<int> z = zFunction(text);
\end{lstlisting}